\section{Thesis structure}
\label{sec:intro:structure}

The remainder of this thesis is organised into five further
chapters and three appendixes.

\paragraph{Chapter~\ref{chapter:background} --- Background and
Related Work.} Opens with a Related Work survey of the three
bodies of prior research the proposed solution builds upon ---
bus operational analytics, visualisation of spatial-temporal
traffic data, and existing HCMC analytical tools --- and then
develops the theoretical foundations the solution rests on: the
four-layer analytical warehouse, the DuckDB + Streamlit + Plotly
+ pydeck stack, the QoS scoring framework of
Srisakaokul~\textit{et al.}~\cite{srisakaokul2023}, the
visualisation primitives of Liu~\cite{liu2022}
(\emph{Gridify}, \emph{GroupSet}, \emph{FuzzyCut}) and the Marey
time--distance diagram of Lin~\&~Zeng~\cite{linzeng2014}, the
visual-storytelling principles of Knaflic~\cite{knaflic2015}, and
the Multiple Coordinated Views interaction model that the
dashboard's linked-brushing invariant implements.

\paragraph{Chapter~\ref{chapter:proposed} --- Proposed Solution.}
Opens with the problem statement, then develops the proposed
solution along five axes: an overview of the platform and its
solution-decision-value-framework prioritisation; the data
architecture (warehouse schema, conformed dimensions, fact and
aggregate views); the exploratory data analysis that scopes every
downstream design decision; the screening of analytical modules
against schema and EDA; and the dashboard specification (Hub
$\rightarrow$ Card $\rightarrow$ Chart hierarchy, six cards,
thirty-one canonical charts, three behavioural invariants, and
the scrollytelling Stories layer). This is the
highest-weighted chapter of the report (Câu~4 of the grading
rubric~\cite{rubric}, thirty points).

\paragraph{Chapter~\ref{chapter:prototype} --- Current Prototype.}
Presents the artefacts produced during Phase~1: the four
exploratory notebooks, the published high-fidelity UI mockup,
the data-pipeline skeleton, and the supporting design artefacts
that the Phase~2 implementation depends on.

\paragraph{Chapter~\ref{chapter:future} --- Future Work.}
Sketches the Phase~2 implementation roadmap (nine phases,
P1--P9, approximately twenty-one working days), the data
activations required to unlock currently-deferred features (full
50-day GPS, school geocoding), the predictive extensions that
follow naturally from the historical foundation
(arrival-time prediction, anomaly detection), and the planned
external integrations and user study. Together with
Appendix~B, this chapter addresses Câu~5 of the grading
rubric~\cite{rubric}.

\paragraph{Chapter~\ref{chapter:conclusion} --- Conclusion.}
Summarises the work, documents the four limitations that frame
its scope, enumerates the contributions of Phase~1, and closes
the report.

\paragraph{Appendixes.} Appendix~A documents the warehouse DDL
(table definitions for \texttt{dim\_route}, \texttt{dim\_date},
\texttt{fact\_boardings}, and \texttt{agg\_route\_hour}).
Appendix~B presents the Phase~2 plan as a phased schedule with
deliverables, dependencies, and expected outcomes per phase.
Appendix~C is the master inventory of every image and diagram
this report needs, with capture instructions for the author.
